\section{Related Work}\label{sec:related}

\subsection*{Rough Sets, Feature Selection, and Exact Relevance}

At the static level, the closest classical neighbor is rough-set reduct theory~\cite{pawlak1982rough,slowinski1995decision}. That literature studies which attributes can be deleted while preserving decision distinctions. Exact static sufficiency can be viewed as a reduct condition for the induced decision table. The present paper is different in scope: it is not a new reduct algorithm, but a structural study of which tractable mechanisms have been formalized and whether they admit a finite characterization.

A separate neighboring literature studies feature selection, variable importance, and model explanation in machine learning and AI, including wrapper and filter methods~\cite{kohavi1997wrappers}, general surveys~\cite{guyon2003introduction}, and attribution methods based on Shapley-style decompositions~\cite{lundberg2017unified}. These works are relevant by motivation rather than by technique. They typically optimize predictive performance, explanatory salience, or approximation quality under data-dependent criteria. Exact relevance certification instead asks for a semantic preservation property: which coordinates are necessary to preserve the optimizer relation exactly.

\subsection*{Backdoors, Tractable Islands, and Dichotomy Programs}

The nearest complexity-theoretic analogue is backdoor tractability for constraint satisfaction~\cite{williams2003backdoors,bessiere2013detecting,ganian2017discovering}. There, a small structural set exposes membership in a tractable class even when the ambient problem is hard. The comparison is direct: exact relevance certification also asks which coordinates matter, and polynomial-time behavior emerges when the source of hardness is restricted by structure.

The broader analogy is the dichotomy tradition for satisfiability and constraint satisfaction, from Schaefer's Boolean dichotomy theorem~\cite{schaefer1978complexity} to the finite-domain CSP dichotomy theorems of Bulatov and Zhuk~\cite{bulatov2017dichotomy,zhuk2017proof}. The expectation that a clean tractability boundary should exist was articulated in the Feder--Vardi program~\cite{feder1998computational}, and later algebraic work of Barto and Kozik clarified why finite structural boundaries are plausible in that setting~\cite{barto2014constraint}. The present paper should be read against that background. It does not argue against dichotomy theorems in general. It shows instead that, for exact relevance certification, quotient realizability and closure-law invariance leave too much room for a naive admissible classifier to survive. Any successful frontier theorem in this domain must therefore use stronger representation-sensitive structure than the first natural candidate provides.

\subsection*{Meta-Impossibility Traditions}

The negative theorem in this paper belongs to a broader family of impossibility results. At the foundational level, G"odel's incompleteness theorems show that sufficiently expressive formal systems cannot internalize all truths about their own proof theory~\cite{goedel1931formal}. In learning and optimization, no-free-lunch results show that universal guarantees disappear on overly broad problem domains~\cite{wolpert1996lack,wolpert1997nfl}. In social choice, Arrow's theorem and the Gibbard--Satterthwaite theorems show that natural axiomatic requirements become jointly unsatisfiable on rich enough domains~\cite{arrow1950difficulty,gibbard1973manipulation,satterthwaite1975strategy}. The present theorem has the same broad form: one fixes a principled invariance package and then proves that no exact classifier can satisfy it on the intended domain.

The closest analogy is nevertheless Rice-style impossibility. Rice's theorem~\cite{rice1953classes} shows that nontrivial extensional properties of sufficiently expressive semantic objects are undecidable. Refinements such as the Myhill--Shepherdson theorem and the Rice--Shapiro theorem show that this is not a single isolated fact but part of a broader pattern for structured semantic domains~\cite{myhill1955effective,shapiro1956degrees}. The analogue in the present paper is that dominant-pair concentration, margin-boundedness, ghost-action concentration, and additive/statewise offset concentration are nontrivial slice predicates that are invariant under the theorem-forced closure laws.

The Arrow parallel is complementary. The admissibility package plays the role of the axiom system: closure-law invariance formalizes semantic irrelevance, while polynomial-time checkability, structural extractability, and bounded-pattern definability exclude semantic encodings that would make the statement vacuous. Once those requirements are imposed, the orbit-gap witnesses show that exact correctness is impossible. The result is therefore best understood not as a pathology of one chosen classifier, but as a structural impossibility theorem for a natural class of representation-sensitive classifiers.
